\chapter{評価}\label{cha:Evaluation}

本研究で提案するプロンプトの有効性を検証するため、被験者を用いた比較実験を行った。
実験では、実際のシステム開発を想定した題材を用い、提案手法とその他の手法で作成された章立ての品質を専門家が評価した。

\section{実験設定}

\subsection{対象システム}
実験の題材として、以下の2つのシステムを選定した。
\begin{enumerate}
    \item \textbf{BMP画像分割アプリ}: BMP画像を指定された数に分割し、保存するアプリケーション
    \item \textbf{クラウド型電子契約サービス}: PDF技術を活かした電子契約／ワークフロー、それらに付随する文書管理、及びユーザー管理機能を持つクラウドサービス
\end{enumerate}

\subsection{比較手法}
作成手法として、以下の3パターンを設定した。
\begin{itemize}
    \item \textbf{手法A（手動作成）}: 被験者はLLMを用いず、章立てを作成する。
    \item \textbf{手法B（LLMのみ）}: 被験者はLLMを使用するが、本研究のプロンプトは使用せず、自由な指示で章立てを作成する。
    \item \textbf{手法C（提案手法）}: 被験者はLLMを使用し、本研究で作成したテンプレートプロンプトを用いて章立てを作成する。
\end{itemize}
なお、被験者CとDの手法Bにおいて、プロンプトの指定はないものとする。

\subsection{被験者と実験の割り当て}
4名の被験者（A, B, C, D）に対し、システムと手法の組み合わせを以下の表\ref{tb:experiment_allocation}のように割り当てた。
クロス集計的な評価が可能になるように、被験者ごとに異なる組み合わせを設定している。

\begin{table}[h]
  \caption{被験者ごとのシステムと採用手法の割り当て}
  \label{tb:experiment_allocation}
  \centering
  \begin{tabular}{|c|l|l|}
    \hline
    \textbf{被験者} & \textbf{BMP画像分割アプリ} & \textbf{クラウド型電子契約サービス} \\
    \hline\hline
    A & 手法A（手動） & 手法C（提案手法） \\
    \hline
    B & 手法C（提案手法） & 手法A（手動） \\
    \hline
    C & 手法B（LLMのみ） & 手法C（提案手法） \\
    \hline
    D & 手法C（提案手法） & 手法B（LLMのみ） \\
    \hline
  \end{tabular}
\end{table}

\section{評価基準}
作成された章立ての成果物は、株式会社スカイコムにて5年以上のシステム仕様書作成経験を持つ専門家によって評価された。
評価は以下の8項目について行われ、それぞれ5段階のリッカート尺度（5: 非常に良い 〜 1: 非常に悪い）で採点された。

\begin{enumerate}
    \item 最低限含むべき章が含まれているか（網羅性）
    \item 章の順序は適切か（論理性）
    \item 章タイトルは適切か（命名の適切さ）
    \item 章立て全体として理解しやすいか（可読性）
    \item 章の粒度は適切か（詳細度）
    \item 想定読者（新人プログラマ等）に適した章立てか（適合性）
    \item 将来的な拡張に対応しやすい章立てか（拡張性）
    \item 章、節、項の包含関係は適切か（構造的整合性）
\end{enumerate}

\section{実験結果}
各手法によって作成されたシステム仕様書章立ての評価結果について、手法別に比較を行う。
まず、手法A（手動作成）と手法C（提案手法）の比較結果を表\ref{tb:evaluation_result_ac}に示す。

\begin{figure}[htbp]
\begin{tcolorbox}[
  colback=white,
  colframe=black,
  fonttitle=\bfseries
]
\begin{multicols}{2}
\footnotesize
\ttfamily
\fontsize{11}{12}\selectfont
\begin{verbatim}
1. 概要
 1.1.目的
 1.2.アプリケーションの概要
2. アプリケーションの基本条件
 2.1.入力画像の条件
 2.2.動作環境
 2.3.性能条件
3. 構成
 1.アプリケーションの構成
 2.ネットワークの構成
4. アプリケーションの機能
 4.1.単一画像分割機能
 4.2.複数枚分割機能
5. アプリケーションの運用
 5.1.運用方法
 5.2.障害時の運用
6.セキュリティ
7.安全設計
8.データ仕様
 8.1.コード体系
 8.2.データベース
 8.3.ファイル
 8.4.入出力データ
 8.5.メッセージ
9.付帯事項
\end{verbatim}
\end{multicols}
\end{tcolorbox}
\caption{手法AによるBMP画像分割アプリの章立て}
\label{fig:A_chapter_structure_of_bmp_app}
\end{figure}

\begin{figure}[htbp]
\begin{tcolorbox}[
  colback=white,
  colframe=black,
  fonttitle=\bfseries
]
\begin{multicols}{2}
\footnotesize
\ttfamily
\fontsize{11}{12}\selectfont
\begin{verbatim}
1. 概要
 1.1 目的
 1.2 システム概要
 1.3 用語・略語定義
2. システム基本条件
 2.1 動作環境
  2.1.1 対応OS
  2.1.2 実行形態
 2.2 対応BMPファイル仕様
  2.2.1 対応ビット深度
  2.2.2 非対応形式の扱い
 2.3 性能条件
  2.3.1 最大画像サイズ
  2.3.2 処理時間目安
 2.4 制約条件
  2.4.1 分割数の制約
  2.4.2 メモリ使用制約
3. システム構成
 3.1 全体構成
  3.1.1 処理フロー
  3.1.2 入出力関係
 3.2 ソフトウェア構成
  3.2.1 使用プログラミング言語
  3.2.2 使用ライブラリ
4. 機能仕様
 4.1 機能一覧
 4.2 画像入力機能
  4.2.1 入力方法
  4.2.2 入力チェック
 4.3 分割数指定機能
  4.3.1 指定方法
  4.3.2 不正値の扱い
 4.4 画像分割処理機能
  4.4.1 分割処理開始条件
  4.4.2 分割処理結果
 4.5 分割画像出力機能
  4.5.1 出力形式
  4.5.2 ファイル命名規則
5. 画像分割処理仕様
 5.1 分割方式
  5.1.1 縦横分割の定義
  5.1.2 均等分割の考え方
 5.2 端数処理
  5.2.1 幅方向の端数処理
  5.2.2 高さ方向の端数処理
 5.3 処理アルゴリズム概要
  5.3.1 ピクセル取得処理
  5.3.2 画像生成処理
6. 入出力仕様
 6.1 入力データ仕様
  6.1.1 入力ファイル
  6.1.2 入力パラメータ
 6.2 出力データ仕様
  6.2.1 出力ファイル構成
  6.2.2 出力先ディレクトリ
7. エラー処理・例外処理
 7.1 入力エラー
  7.1.1 ファイル未存在
  7.1.2 形式不正
 7.2 処理エラー
  7.2.1 メモリ不足
  7.2.2 分割処理失敗
8. 操作・運用
 8.1 操作手順
  8.1.1 起動
  8.1.2 処理実行
  8.1.3 終了
 8.2 エラー発生時の対応
  8.2.1 エラーメッセージ表示
  8.2.2 ユーザー対応
9. セキュリティ・安全設計
 9.1 ファイル操作上の安全性
 9.2 不正入力対策
10. 付帯事項
 10.1 既知の制限事項
 10.2 今後の拡張予定
\end{verbatim}
\end{multicols}
\end{tcolorbox}
\caption{手法BによるBMP画像分割アプリの章立て}
\label{fig:B_chapter_structure_of_bmp_app}
\end{figure}

\begin{figure}[htbp]
\begin{tcolorbox}[
  colback=white,
  colframe=black,
  fonttitle=\bfseries
]
\begin{multicols}{2}
\footnotesize
\ttfamily
\fontsize{11}{12}\selectfont
\begin{verbatim}
1. 概要
　1.1. 目的
　1.2. システム概要
　1.3. 用語定義
　1.4. 適用範囲
2. システム全体像
　2.1. システム構成図
　2.2. 利用者とユースケース
3. 動作環境・前提条件
　3.1. 動作環境
　3.2. 対象ファイル条件
　3.3. 制約事項
4. 画面仕様
　4.1. 画面一覧
　4.2. メイン画面
　　4.2.1. 画面レイアウト
　　4.2.2. 入力項目仕様
　　4.2.3. 操作手順
5. 機能仕様
　5.1. 機能一覧
　5.2. BMP画像分割機能
　　5.2.1. 処理概要
　　5.2.2. 入力仕様
　　5.2.3. 出力仕様
6. エラー仕様
　6.1. エラー一覧
　6.2. 入力値エラー
　6.3. ファイルエラー
　6.4. システムエラー
7. 非機能要件
　7.1. パフォーマンス要件
　7.2. セキュリティ要件
　7.3. 操作性要件
8. データ仕様
　8.1. 入力データ形式
　8.2. 出力データ形式
　8.3. ファイル命名規則
9. 運用仕様
9.1. 通常運用
　9.2. エラー発生時の対応
10. テスト観点整理
　10.1. テスト対象範囲
　10.2. エラー条件と期待結果
\end{verbatim}
\end{multicols}
\end{tcolorbox}
\caption{手法CによるBMP画像分割アプリの章立て1}
\label{fig:C_chapter_structure_of_bmp_app_1}
\end{figure}



\begin{figure}[htbp]
\begin{tcolorbox}[
  colback=white,
  colframe=black,
  fonttitle=\bfseries
]
\begin{multicols}{2}
\footnotesize
\ttfamily
\fontsize{11}{12}\selectfont
\begin{verbatim}
1. 概要
 1.1. 目的
 1.2. 用語定義
 1.3. システム概要
 1.4. 適用範囲
2. 動作環境・前提条件
 2.1. 動作環境
 2.2. 対象ファイル条件
 2.3. 利用上の前提・制約
3. システム全体仕様
 3.1. システム構成概要
 3.2. 処理概要
4. ユースケース
 4.1. ユースケース一覧
 4.2. ユースケース詳細
5. 画面仕様
 5.1. 画面一覧
 5.2. 画面レイアウト
 5.3. 操作フロー
6. 機能仕様
 6.1. 機能一覧
 6.2. 画像分割機能
 6.3. 入力チェック機能
7. 入出力仕様
 7.1. 入力データ仕様
 7.2. 出力データ仕様
 7.3. ファイル命名規則
8. エラー・例外仕様
 8.1. エラー一覧
 8.2. エラーメッセージ仕様
 8.3. エラー発生時の画面・動作
9. 非機能要件
 9.1. 性能要件
 9.2. セキュリティ要件
 9.3. ユーザビリティ要件
10. 運用・テスト観点
 10.1. 想定利用方法
 10.2. テスト観点一覧
11. 付録
 11.1. 参考資料
 11.2. 変更履歴
\end{verbatim}
\end{multicols}
\end{tcolorbox}
\caption{手法CによるBMP画像分割アプリの章立て2}
\label{fig:C_chapter_structure_of_bmp_app_2}
\end{figure}



\begin{figure}[htbp]
\begin{tcolorbox}[
  colback=white,
  colframe=black,
  fonttitle=\bfseries
]
\begin{multicols}{2}
\footnotesize
\ttfamily
\fontsize{11}{12}\selectfont
\begin{verbatim}
1. 概要
 1.1.目的
 1.2.システム概要
 1.3.適用範囲
2. システム基本条件
 2.1.稼働条件
 2.2.他システムインターフェース
 2.3.性能条件
 2.4.拡張条件
 2.5.継承条件
 2.6.以降環境
3. システム構成
 3.1.ハードウェア構成及びネットワーク構成
 3.2.ソフトウェア構成
4. システムの機能
 4.1.機能構成
 4.2.個別機能概要
   4.2.1電子契約機能
   4.2.2ワークフロー生成機能
   4.2.3文書管理機能
   4.2.4ユーザ管理機能
5. システムの運用
 5.1.運用方法
 5.2.各業務(機能)の運用
 5.3.障害時の運用
6.セキュリティ
7.安全設計
8.データ仕様
 8.1.コード体系
 8.2.データベース
 8.3.ファイル
 8.4.入出力データ
 8.5.メッセージ
9.付帯事項
\end{verbatim}
\end{multicols}
\end{tcolorbox}
\caption{手法Aによるクラウド型電子契約システムの章立て}
\label{fig:A_chapter_structure_of_cloud_contract_system}
\end{figure}



\begin{figure}[htbp]
\begin{tcolorbox}[
  colback=white,
  colframe=black,
  fonttitle=\bfseries
]
\begin{multicols}{2}
\footnotesize
\ttfamily
\fontsize{11}{12}\selectfont
\begin{verbatim}
1. 概要
 1.1 目的
 1.2 システム概要
 1.3 用語定義
 1.4 適用範囲
 1.5 想定利用者
2. システム基本条件
 2.1 稼働環境・前提条件
 2.2 対応ブラウザ・デバイス
 2.3 他システム連携（外部インターフェース）
   2.3.1 電子署名・認証基盤
   2.3.2 メール／通知サービス
   2.3.3 外部API連携
 2.4 性能・可用性要件
 2.5 拡張性・将来要件
 2.6 移行・継承条件
3. システム構成
 3.1 全体構成図
 3.2 クラウド基盤構成
 3.3 ネットワーク構成
 3.4 ソフトウェア構成
 3.5 利用サービス・ミドルウェア一覧
4. 機能仕様
 4.1 機能一覧
 4.2 契約作成・アップロード機能
 4.3 電子署名・押印機能
 4.4 ワークフロー管理機能
 4.5 契約締結・完了処理
 4.6 文書管理機能
 4.7 ユーザー・権限管理機能
 4.8 通知・リマインド機能
 4.9 検索・閲覧機能
 4.10 管理者機能
5. 画面仕様
 5.1 画面一覧
 5.2 画面遷移図
 5.3 各画面詳細
   5.3.1 契約一覧画面
   5.3.2 契約詳細画面
   5.3.3 署名画面
   5.3.4 管理画面
6. データ仕様
 6.1 データ構造概要
 6.2 データベース設計
 6.3 契約データ仕様
 6.4 ユーザーデータ仕様
 6.5 ファイル仕様（PDF等）
 6.6 入出力データ定義
7. セキュリティ設計
 7.1 認証・認可方式
 7.2 電子署名の真正性・完全性
 7.3 通信・保存時の暗号化
 7.4 アクセス制御
 7.5 監査ログ・操作ログ
8. 安全性・法的要件
 8.1 電子帳簿保存法対応
 8.2 電子署名法対応
 8.3 タイムスタンプ要件
 8.4 データ保全・バックアップ
9. 運用・保守
 9.1 運用体制
 9.2 日常運用
 9.3 障害対応
 9.4 バックアップ・リストア
 9.5 監視・ログ管理
10. 非機能要件
 10.1 性能要件
 10.2 可用性要件
 10.3 保守性要件
 10.4 信頼性要件
11. 付帯事項
 11.1 制約事項
 11.2 未決事項
 11.3 参考資料
\end{verbatim}
\end{multicols}
\end{tcolorbox}
\caption{手法Bによるクラウド型電子契約システムの章立て}
\label{fig:B_chapter_structure_of_cloud_contract_system}
\end{figure}



\begin{figure}[htbp]
\begin{tcolorbox}[
  colback=white,
  colframe=black,
  fonttitle=\bfseries
]
\begin{multicols}{2}
\footnotesize
\ttfamily
\fontsize{11}{12}\selectfont
\begin{verbatim}
1. 概要
 1.1. 目的
 1.2. 用語定義
 1.3. システム概要
 1.4. 適用範囲
 1.5. 前提条件・制約条件
2. システム全体像
 2.1. システム構成図
 2.2. 利用者区分
 2.3. ユースケース一覧
 2.4. ユースケース図
3. 動作環境・基本条件
 3.1. 動作環境
 3.2. 稼働時間・可用性
 3.3. 性能要件
 3.4. 拡張性・将来対応
4. 画面仕様
 4.1. 画面一覧
 4.2. 共通画面仕様
 4.3. 個別画面仕様
  4.3.1. ログイン画面
  4.3.2. 契約書アップロード画面
  4.3.3. 電子署名画面
  4.3.4. ワークフロー管理画面
  4.3.5. 文書一覧・検索画面
  4.3.6. ユーザー管理画面
5. 機能仕様
 5.1. 機能構成
 5.2. 電子契約機能
 5.3. ワークフロー機能
 5.4. 文書管理機能
 5.5. ユーザー管理機能
 5.6. 通知機能
6. 操作フロー仕様
 6.1. 契約締結フロー
 6.2. 承認・差戻しフロー
 6.3. 再署名・再送フロー
7. 外部インターフェース仕様
 7.1. 外部システム連携概要
 7.2. API仕様
 7.3. ファイル入出力仕様
8. データ仕様
 8.1. データ項目定義
 8.2. コード体系
 8.3. 入出力データ定義
9. エラー・メッセージ仕様
 9.1. エラー種別
 9.2. エラーメッセージ一覧
 9.3. エラー時の画面動作
10. セキュリティ仕様
 10.1. 認証・認可
 10.2. データ保護
 10.3. 操作ログ・監査ログ
11. 運用仕様
 11.1. 通常運用
 11.2. 障害時運用
 11.3. メンテナンス運用
12. テスト観点整理
 12.1. 機能テスト観点
 12.2. 異常系テスト観点
13. 付録
 13.1. 画面遷移一覧
 13.2. 用語一覧
\end{verbatim}
\end{multicols}
\end{tcolorbox}
\caption{手法Cによるクラウド型電子契約システムの章立て1}
\label{fig:C_chapter_structure_of_cloud_contract_system_1}
\end{figure}


\begin{figure}[htbp]
\begin{tcolorbox}[
  colback=white,
  colframe=black,
  fonttitle=\bfseries
]
\begin{multicols}{2}
\footnotesize
\ttfamily
\fontsize{11}{12}\selectfont
\begin{verbatim}
1. 概要
　1.1. 目的
　1.2. システム概要
　1.3. 用語定義
　1.4. 適用範囲
2. システム全体像
　2.1. システム構成概要
　2.2. ユースケース一覧
　2.3. 利用者区分
3. システム基本条件
　3.1. 動作環境
　3.2. ネットワーク条件
　3.3. 他システム連携条件
　3.4. 制約事項
4. 機能構成
　4.1. 機能一覧
　4.2. 機能間関係
5. 個別機能仕様
　5.1. 電子契約機能
 　5.1.1. 契約書作成
 　5.1.2. 署名・押印
 　5.1.3. 契約締結
　5.2. ワークフロー機能
 　5.2.1. 承認フロー定義
 　5.2.2. 承認・却下操作
　5.3. 文書管理機能
 　5.3.1. 文書保管
 　5.3.2. 文書検索
　5.4. ユーザー管理機能
 　5.4.1. ユーザー登録
 　5.4.2. 権限管理
6. 画面仕様
　6.1. 画面一覧
　6.2. 画面共通仕様
　6.3. 画面別レイアウト
7. 操作フロー
　7.1. 通常操作フロー
　7.2. 異常時操作フロー
8. データ仕様
　8.1. データ項目定義
　8.2. ファイル仕様
　8.3. 入出力データ仕様
9. エラー仕様
　9.1. エラー一覧
　9.2. エラーメッセージ
10. 非機能要件
　10.1. パフォーマンス要件
　10.2. セキュリティ要件
　10.3. 可用性要件
11. 運用仕様
　11.1. 運用方法
　11.2. 障害時対応
12. テスト観点
　12.1. 機能テスト観点
　12.2. 異常系テスト観点
13. 付録
　13.1. 参考資料
　13.2. 変更履歴
\end{verbatim}
\end{multicols}
\end{tcolorbox}
\caption{手法Cによるクラウド型電子契約システムの章立て2}
\label{fig:C_chapter_structure_of_cloud_contract_system_2}
\end{figure}


\begin{table}[h]
  \caption{手法A（手動）と手法C（提案手法）の平均スコア比較（ダミーデータ）}
  \label{tb:evaluation_result_ac}
  \centering
  \begin{tabular}{|l|c|c|c|}
    \hline
    \textbf{評価項目} & \textbf{手法A (手動)} & \textbf{手法C (提案手法)} & \textbf{差 (C - A)} \\
    \hline\hline
    1. 網羅性 & 3.5 & \textbf{4.8} & +1.3 \\
    \hline
    2. 論理性 & 4.0 & \textbf{4.5} & +0.5 \\
    \hline
    3. 命名の適切さ & 3.8 & \textbf{4.2} & +0.4 \\
    \hline
    4. 可読性 & 3.2 & \textbf{4.6} & +1.4 \\
    \hline
    5. 詳細度 & 3.0 & \textbf{4.0} & +1.0 \\
    \hline
    6. 適合性 & 3.5 & \textbf{4.7} & +1.2 \\
    \hline
    7. 拡張性 & 3.0 & \textbf{4.3} & +1.3 \\
    \hline
    8. 構造的整合性 & 3.8 & \textbf{4.9} & +1.1 \\
    \hline
  \end{tabular}
\end{table}

次に、手法B（LLMのみ）と手法C（提案手法）の比較結果を表\ref{tb:evaluation_result_bc}に示す。

\begin{table}[h]
  \caption{手法B（LLMのみ）と手法C（提案手法）の平均スコア比較（ダミーデータ）}
  \label{tb:evaluation_result_bc}
  \centering
  \begin{tabular}{|l|c|c|c|}
    \hline
    \textbf{評価項目} & \textbf{手法B (LLMのみ)} & \textbf{手法C (提案手法)} & \textbf{差 (C - B)} \\
    \hline\hline
    1. 網羅性 & 3.8 & \textbf{4.8} & +1.0 \\
    \hline
    2. 論理性 & 3.0 & \textbf{4.5} & +1.5 \\
    \hline
    3. 命名の適切さ & 3.5 & \textbf{4.2} & +0.7 \\
    \hline
    4. 可読性 & 3.5 & \textbf{4.6} & +1.1 \\
    \hline
    5. 詳細度 & 2.5 & \textbf{4.0} & +1.5 \\
    \hline
    6. 適合性 & 3.2 & \textbf{4.7} & +1.5 \\
    \hline
    7. 拡張性 & 3.5 & \textbf{4.3} & +0.8 \\
    \hline
    8. 構造的整合性 & 3.0 & \textbf{4.9} & +1.9 \\
    \hline
  \end{tabular}
\end{table}




図\ref{fig:result_chart}に、各手法の評価結果をレーダーチャートで比較したものを示す。

\begin{figure}[htbp]
  \centering
  % TODO: ここにExcel等で作成したレーダーチャートの画像を差し替える (images/result_chart.pdf)
  \includegraphics[width=0.8\linewidth]{images/syokuji_computer.png} % 仮の画像
  \caption{実験結果の比較（レーダーチャート）}
  \label{fig:result_chart}
\end{figure}

実験の結果、提案手法（手法C）は、すべての評価項目において手法Aおよび手法Bを上回る、または同等のスコアを記録した。
特に「網羅性」と「構造的整合性」において顕著な高さを示しており、これはプロンプトで与えた「ベース章立て」と「出力指示」の有効性を示唆している。
一方で、「詳細度」に関しては手法A（手動）との差が比較的小さく、人間のドメイン知識には及ばない部分があることも確認された。
